\chapter{Bacterial Vaginosis}

\section*{What Is This Test?}
Bacterial vaginosis (BV) testing assesses the shift in the vaginal microbiome from a \textit{Lactobacillus}-dominant environment (healthy vaginal flora) to a polymicrobial dysbiosis characterized by a marked reduction in hydrogen peroxide-producing \textit{Lactobacillus} species (especially \textit{L. crispatus, L. jensenii, L. gasseri}) and a 100- to 1,000-fold overgrowth of facultative and obligate anaerobic bacteria, including \textit{Gardnerella vaginalis} (the most prominent but not sufficient cause), \textit{Atopobium vaginae}, \textit{Prevotella} species, \textit{Mobiluncus} species, \textit{Bacteroides} species, \textit{Megasphaera}, \textit{Peptostreptococcus}, \textit{Sneathia}, \textit{BVAB1-3}, and \textit{Mycoplasma hominis}. BV is the most common cause of vaginal discharge in women of reproductive age, affecting approximately 30\% of women in the US (higher in Black and Hispanic women). It is not a classic STI (it can occur in women who have never been sexually active), but sexual activity (new partner, multiple partners, lack of condom use, receptive oral sex) is a major risk factor. BV is associated with adverse outcomes including preterm birth, premature rupture of membranes, postpartum endometritis, post-hysterectomy vaginal cuff cellulitis, and increased risk of acquiring HIV, HSV-2, chlamydia, gonorrhea, and trichomoniasis.

\section*{Alternative Names}
BV Test; Gram Stain for Vaginosis; Nugent Score; Amsel Criteria; Vaginal pH; Whiff Test; Clue Cells; BV NAAT; Bacterial Vaginosis PCR; Vaginitis Panel; OSOM BVBlue Test

\section*{Principle}
Diagnosis is established by clinical criteria (Amsel) or laboratory testing (Gram stain with Nugent score; NAAT). Amsel clinical criteria: Requires 3 of the 4 following criteria to be present: (1) Thin, homogeneous, white-gray, non-inflammatory vaginal discharge that smoothly coats the vaginal walls; (2) Vaginal pH >4.5 (measured by pH paper applied to the lateral vaginal wall, not the cervical os or posterior fornix where cervical mucus/ semen/ blood can falsely elevate pH); (3) Positive whiff (amine) test: release of a fishy (amine) odor upon addition of 10\% KOH to a sample of vaginal discharge; (4) Clue cells on saline wet mount: vaginal squamous epithelial cells with edges obscured by adherent coccobacilli (comprising >20\% of epithelial cells). Sensitivity 35--90\%, specificity 60--90\% depending on the clinician's experience. Gram stain with Nugent scoring: Vaginal smear Gram stained, and bacterial morphotypes are quantified under oil immersion (1,000x). The Nugent score (0--10) is calculated as the sum of scores for: (a) Large gram-positive rods (\textit{Lactobacillus} morphotypes): 0--4+; (b) Small gram-variable or gram-negative coccobacilli (\textit{Gardnerella/Bacteroides} morphotypes): 0--4+; (c) Curved gram-variable or gram-negative rods (\textit{Mobiluncus} morphotypes): 0--2+. Score 0--3: Normal flora; 4--6: Intermediate; 7--10: BV. Nugent score is the gold standard for research and more objective than Amsel criteria. Sensitivity 89\%, specificity 83\%. NAAT: Multiplex PCR panels (BD MAX Vaginal Panel, Hologic Aptima BV, Cepheid Xpert Xpress MVP, BioFire FilmArray Vaginitis Panel) detect and quantitate \textit{Lactobacillus} species and BV-associated bacteria (e.g., \textit{Gardnerella vaginalis}, \textit{Atopobium vaginae}, \textit{Megasphaera} type 1, BVAB2). These assays report a qualitative or semi-quantitative BV result. FDA-cleared NAATs have sensitivity 90--97\% and specificity 85--95\%. The Affirm VPIII DNA probe detects \textit{G. vaginalis} (at high concentrations >2 x 10$^5$ CFU/mL) but is inferior to NAAT and is being phased out. BVBlue (OSOM): Detects sialidase activity, an enzyme produced by BV-associated anaerobes; results in 10 minutes with sensitivity 88\% and specificity 95\%.

\section*{History}
Bacterial vaginosis was originally termed "nonspecific vaginitis" and was attributed to \textit{Haemophilus vaginalis} (later renamed \textit{Gardnerella vaginalis}) by Herman Gardner in 1955 \cite{gardner1955haemophilus}. It was renamed "bacterial vaginosis" in 1984 to acknowledge that it is a dysbiosis (not a single infection) and that inflammation is absent (vagin-osis, not vagin-itis). The Amsel clinical criteria were published in 1983 \cite{amsel1983nonspecific}. The Nugent scoring system was published in 1991 \cite{nugent1991reliability}, providing a more objective, reproducible gold standard for diagnosis. The recognition of \textit{Atopobium vaginae} and other fastidious anaerobes as key players in BV pathogenesis came with molecular techniques (16S rRNA gene sequencing) in the 2000s. Multiplex molecular panels for vaginitis (including BV, \textit{Candida} species, and \textit{T. vaginalis}) received FDA clearance in the 2010s.

\section*{Why This Test Is Ordered}
\begin{itemize}
  \item Evaluation of women complaining of a malodorous, thin, homogeneous, gray-white vaginal discharge (often described as a "fishy" odor), particularly noted after unprotected intercourse (semen alkalinizes the vagina, volatilizing amines)
  \item Investigation of persistent or recurrent vaginal discharge with negative testing for \textit{Candida} (yeast) and \textit{Trichomonas vaginalis}
  \item Preoperative screening prior to hysterectomy (total abdominal, vaginal, or laparoscopic) or pregnancy termination, as BV increases the risk of post-procedural pelvic infection (vaginal cuff cellulitis, PID) 3- to 5-fold. The ACOG recommends screening and treatment of BV before hysterectomy and elective pregnancy termination
  \item Evaluation of pregnant women with a history of preterm birth; screening for BV in the early second trimester and treating if positive has been shown in some, but not all, studies to reduce the risk of preterm birth in high-risk women, though universal screening of asymptomatic pregnant women is not recommended by USPSTF or ACOG
  \item Evaluation of women with HIV infection and vaginal symptoms, as BV is more prevalent and persistent in HIV-infected women and may increase HIV genital tract shedding
\end{itemize}

\section*{Primary vs Secondary Test}
The clinical Amsel criteria, Gram stain with Nugent score, and NAAT are all primary diagnostic methods for BV. NAAT offers the highest sensitivity and is increasingly used in clinical laboratories. Wet-mount examination for clue cells, pH measurement, and whiff test are part of the clinical "bedside" test that constitutes the Amsel criteria and requires no laboratory infrastructure beyond a microscope and pH paper.

\section*{How This Test Is Performed}
The patient is placed in the dorsal lithotomy position. A sterile, non-lubricated speculum (lubricant can be bactericidal and may interfere with culture and pH) is inserted. The character of the vaginal discharge is noted (color, consistency, adherence to the vaginal walls). A sterile cotton- or Dacron-tipped swab is used to collect a sample of vaginal discharge from the lateral vaginal wall and the posterior fornix. The swab is used to: (1) roll onto a glass slide for Gram stain (which can be used for Nugent scoring); (2) mix with a drop of normal saline on a second slide for wet-mount microscopic examination (for clue cells, trichomonads, and yeast); (3) apply to pH paper (pH is measured by pressing pH paper against the lateral vaginal wall, not the posterior fornix or cervix); (4) perform a whiff test: a drop of 10\% KOH is added to the vaginal discharge on the slide or swab; the immediate release of a fishy, amine-like odor is a positive whiff test. If NAAT is ordered, the swab is placed in the manufacturer-specific transport medium. A single swab is often sufficient for all components if sampled adequately. Self-collected vaginal swabs for NAAT have equivalent sensitivity to clinician-collected swabs.

\section*{What Preparation Is Needed}
The patient should avoid douching, using vaginal creams, suppositories, spermicides, or engaging in sexual intercourse for 24--48 hours before specimen collection. The specimen should ideally not be collected during menstruation (blood interferes with pH measurement, the whiff test, and microscopic examination). The speculum should be inserted without lubricant; if lubrication is necessary, water only or a small amount of non-bacteriostatic lubricant applied to the posterior blade only may be used.

\section*{Contraindications and Precautions}
A positive whiff test can also be produced when blood (alkaline) or semen (alkaline) is present. Vaginal pH is elevated during menses, after intercourse (<24 hours), in the presence of semen, in atrophic vaginitis (postmenopausal, due to reduced \textit{Lactobacillus}), and in trichomoniasis. Elevated pH alone is the least specific Amsel criterion and should not be used in isolation. The presence of trichomonads on wet mount strongly suggests concurrent \textit{T. vaginalis} infection (60--80\% of women with trichomoniasis also have BV); treatment for both is required. BV is frequently associated with an elevated number of WBCs on wet mount, which technically violates the "non-inflammatory" definition of BV: concurrent cervicitis, trichomoniasis, or desquamative inflammatory vaginitis may be present. Recurrent BV (3 or more episodes in 12 months) is common due to persistence of the BV-associated biofilm on the vaginal epithelium despite treatment; relapse rates approach 50--60\% within 6--12 months. The role of \textit{Gardnerella vaginalis}-dominant biofilms adherent to the vaginal epithelium contributes to treatment failure despite \textit{in-vitro} metronidazole susceptibility of the planktonic organisms.

\section*{Risks}
Vaginal speculum examination and swab collection are not painful for most women, though some discomfort may occur. There is no risk from the laboratory tests. The major clinical risk is the failure to diagnose and treat BV before gynecologic surgery, as post-operative infectious morbidity is significantly increased in women with untreated BV.

\section*{What Happens After the Test}
Amsel criteria results are available immediately at the bedside. Gram stain with Nugent score: 24--48 hours. NAAT results: 1--4 hours. Confirmed BV (by Amsel, Nugent, or NAAT) should be treated. CDC-recommended regimens: Metronidazole 500 mg PO BID for 7 days (preferred), OR metronidazole gel 0.75\% (5 g) intravaginally once daily for 5 days, OR clindamycin cream 2\% (5 g) intravaginally at bedtime for 7 days. Alternative: clindamycin 300 mg PO BID for 7 days, OR clindamycin ovules 100 mg intravaginally at bedtime for 3 days, OR secnidazole 2 g PO as a single dose, OR tinidazole 2 g PO daily for 2 days, OR tinidazole 1 g PO daily for 5 days. Alcohol should be avoided during systemic metronidazole/tinidazole therapy and for 24--72 hours afterward. Clindamycin cream and ovules are oil-based and may weaken latex condoms and diaphragms (for 5 days after use for cream, 72 hours for ovules). Treatment of male sexual partners does NOT reduce BV recurrence and is NOT recommended (well-established from multiple RCTs). However, female sexual partners of women with BV are at high risk for BV and should be evaluated.

\section*{Understanding Results}

\subsection*{Below Normal}
\begin{itemize}
  \item \textbf{Amsel criteria:} 0--2 of 4 criteria met. \textbf{Nugent score 0--3:} Normal vaginal flora. \textbf{NAAT negative for BV:} No evidence of BV. These indicate a \textit{Lactobacillus}-predominant healthy vaginal microbiome. If the patient has discharge, consider other etiologies: vulvovaginal candidiasis (yeast), trichomoniasis, desquamative inflammatory vaginitis, atrophic vaginitis, or physiologic discharge
  \item \textbf{Nugent score 4--6 (intermediate):} Not normal and not fully BV. Represents a transitional state that may progress to or regress from BV. In symptomatic women, treatment may be considered. In pregnant women with prior preterm birth, some guidelines recommend treatment
\end{itemize}

\subsection*{Above Normal}
\begin{itemize}
  \item \textbf{Amsel criteria: 3 or 4 of 4 criteria met.} Diagnostic of BV. The presence of clue cells (>20\% of epithelial cells) is the single most reliable criterion; pH >4.5 is the least specific. The whiff test is positive in ~90\% of BV cases
  \item \textbf{Nugent score 7--10:} Diagnostic of BV. The Nugent score is the gold standard for epidemiologic studies and clinical trials and correlates well with BV symptoms
  \item \textbf{NAAT positive for BV (e.g., Aptima BV, BD MAX):} Diagnostic of BV. NAAT is more sensitive than Amsel criteria and eliminates the subjectivity of microscopic examination. The molecular assays detect specific BV-associated bacterial species at concentrations above diagnostic thresholds. The BV panel of the BioFire FilmArray Vaginitis Panel reports \textit{Lactobacillus} species depletion with \textit{G. vaginalis} and other BV-associated anaerobe overgrowth; it also simultaneously tests for \textit{Candida} species and \textit{T. vaginalis} on the same cartridge
  \item \textbf{BVBlue (sialidase) positive:} Consistent with BV; sensitivity comparable to Amsel criteria for symptomatic women but limited sensitivity in asymptomatic women
\end{itemize}

\section*{Normal Results}
Amsel criteria: \textbf{0 of 4 criteria met} (or $\leq$ 2 criteria not diagnostic). Nugent score: \textbf{0--3} (normal \textit{Lactobacillus}-dominant flora). NAAT: \textbf{BV not detected} or \textbf{Normal vaginal flora}. Vaginal pH: \textbf{3.8--4.5}. Whiff test: \textbf{Negative}. Clue cells: \textbf{Absent} (<1\% of epithelial cells).

\section*{Follow-up Tests}
\begin{itemize}
  \item \textbf{Nugent score Gram stain:} If the diagnosis is uncertain by Amsel criteria, a vaginal smear for Nugent scoring provides objective confirmation; this is the preferred method for research settings and for tracking response to therapy
  \item \textbf{Multiplex vaginitis panel (NAAT):} If performed, the panel simultaneously detects BV, \textit{Candida} species, and \textit{T. vaginalis}, eliminating the need for multiple separate tests and providing a comprehensive evaluation for the three most common causes of vaginitis
  \item \textit{\textbf{Trichomonas vaginalis}} \textbf{NAAT:} Trichomoniasis must be excluded, as it commonly co-exists with BV (60--80\% co-infection) and produces similar symptoms (frothy, malodorous discharge) and an elevated pH; treatment for both is required
  \item \textit{\textbf{Candida}} \textbf{species testing (KOH prep, fungal culture, or NAAT):} To evaluate for concurrent vulvovaginal candidiasis, especially if the patient has pruritus and thick, curd-like discharge
  \item \textbf{HIV, syphilis, chlamydia, gonorrhea testing:} BV is associated with increased risk of STI acquisition; women diagnosed with BV, particularly those with risk factors, should be offered comprehensive STI screening
\end{itemize}

\section*{References}
\bibliographystyle{unsrtnat}
\bibliography{references}

\clearpage
